% GENERATED FILE. Edit canonical lessons, then run npm run generate:books.
% canonical-source-hash: fnv1a64:c17c7a30a9fc36a4
% canonical-lessons: SA-C122-rajatam, SA-C122-kacah, SA-C122-sakatam, SA-C122-bhandam, SA-C122-khatva

\chapter{Silver, Glass, Cart, Pot, Cot}
\label{ch:sa-silver-glass-cart-pot-cot}
\hlchaptermodality{122}

\begin{chapteropening}
\textbf{By the end of this chapter:} \emph{I can say silver, glass, a cart, a pot and a cot.}

Complete the last lesson of chapter 122: i can say silver, glass, a cart, a pot and a cot.
\end{chapteropening}

\section[\texorpdfstring{\sk{रजतम्} (rajatam)}{रजतम् (rajatam)}]{\sk{रजतम्} (rajatam) — silver}

\label{lesson:SA-C122-rajatam}



\begin{quote}
Before the new one: say the Sanskrit for a jewel, then the Sanskrit for gold.
\end{quote}

\subsection*{You'll want to know: \sk{रजतम्}}
\textbf{\sk{रजतम्}} — \emph{rajatam} — \textquotedblleft{}silver\textquotedblright{}.

\begin{grammarlens}[title={What you've built}]
One more word for this chapter.
\end{grammarlens}

\subsection*{Your turn}
\begin{itemize}
\raggedright
  \item \emph{Say it:} \emph{rajatam}
  \item \emph{Say it:} \emph{rajatam}, once more
  \item \emph{Say it:} say \emph{rajatam}
  \item \emph{Recall:} say the Sanskrit for a jewel, then the Sanskrit for gold, then say \emph{rajatam} again
\end{itemize}

\begin{tcolorbox}[breakable,colback=teal!4,colframe=teal!35!black,title={Before you move on}]
What does \sk{रजतम्} mean? (\textquotedblleft{}Silver\textquotedblright{}.) Say it once more.
\end{tcolorbox}
\section[\texorpdfstring{\sk{काचः} (kācaḥ)}{काचः (kācaḥ)}]{\sk{काचः} (kācaḥ) — glass}

\label{lesson:SA-C122-kacah}



\begin{quote}
Before the new one: say the Sanskrit for gold, then the Sanskrit for silver.
\end{quote}

\subsection*{You'll want to know: \sk{काचः}}
\textbf{\sk{काचः}} — \emph{kācaḥ} — \textquotedblleft{}glass\textquotedblright{}.

\begin{grammarlens}[title={What you've built}]
One more word for this chapter.
\end{grammarlens}

\subsection*{Your turn}
\begin{itemize}
\raggedright
  \item \emph{Say it:} \emph{kācaḥ}
  \item \emph{Say it:} \emph{kācaḥ}, once more
  \item \emph{Say it:} say \emph{kācaḥ}
  \item \emph{Recall:} say the Sanskrit for gold, then the Sanskrit for silver, then say \emph{kācaḥ} again
\end{itemize}

\begin{tcolorbox}[breakable,colback=teal!4,colframe=teal!35!black,title={Before you move on}]
What does \sk{काचः} mean? (\textquotedblleft{}Glass\textquotedblright{}.) Say it once more.
\end{tcolorbox}
\section[\texorpdfstring{\sk{शकटम्} (śakaṭam)}{शकटम् (śakaṭam)}]{\sk{शकटम्} (śakaṭam) — a cart}

\label{lesson:SA-C122-sakatam}



\begin{quote}
Before the new one: say the Sanskrit for silver, then the Sanskrit for glass.
\end{quote}

\subsection*{You'll want to know: \sk{शकटम्}}
\textbf{\sk{शकटम्}} — \emph{śakaṭam} — \textquotedblleft{}a cart\textquotedblright{}.

\begin{grammarlens}[title={What you've built}]
One more word for this chapter.
\end{grammarlens}

\subsection*{Your turn}
\begin{itemize}
\raggedright
  \item \emph{Say it:} \emph{śakaṭam}
  \item \emph{Say it:} \emph{śakaṭam}, once more
  \item \emph{Say it:} say \emph{śakaṭam}
  \item \emph{Recall:} say the Sanskrit for silver, then the Sanskrit for glass, then say \emph{śakaṭam} again
\end{itemize}

\begin{tcolorbox}[breakable,colback=teal!4,colframe=teal!35!black,title={Before you move on}]
What does \sk{शकटम्} mean? (\textquotedblleft{}A cart\textquotedblright{}.) Say it once more.
\end{tcolorbox}
\section[\texorpdfstring{\sk{भाण्डम्} (bhāṇḍam)}{भाण्डम् (bhāṇḍam)}]{\sk{भाण्डम्} (bhāṇḍam) — a pot}

\label{lesson:SA-C122-bhandam}



\begin{quote}
Before the new one: say the Sanskrit for glass, then the Sanskrit for a cart.
\end{quote}

\subsection*{You'll want to know: \sk{भाण्डम्}}
\textbf{\sk{भाण्डम्}} — \emph{bhāṇḍam} — \textquotedblleft{}a pot\textquotedblright{}.

\begin{grammarlens}[title={What you've built}]
One more word for this chapter.
\end{grammarlens}

\subsection*{Your turn}
\begin{itemize}
\raggedright
  \item \emph{Say it:} \emph{bhāṇḍam}
  \item \emph{Say it:} \emph{bhāṇḍam}, once more
  \item \emph{Say it:} say \emph{bhāṇḍam}
  \item \emph{Recall:} say the Sanskrit for glass, then the Sanskrit for a cart, then say \emph{bhāṇḍam} again
\end{itemize}

\begin{tcolorbox}[breakable,colback=teal!4,colframe=teal!35!black,title={Before you move on}]
What does \sk{भाण्डम्} mean? (\textquotedblleft{}A pot\textquotedblright{}.) Say it once more.
\end{tcolorbox}
\section[\texorpdfstring{\sk{खट्वा} (khaṭvā)}{खट्वा (khaṭvā)}]{\sk{खट्वा} (khaṭvā) — a cot}

\label{lesson:SA-C122-khatva}



\begin{quote}
Before the new one: say the Sanskrit for a cart, then the Sanskrit for a pot.
\end{quote}

\subsection*{You'll want to know: \sk{खट्वा}}
\textbf{\sk{खट्वा}} — \emph{khaṭvā} — \textquotedblleft{}a cot\textquotedblright{}.

\begin{grammarlens}[title={What you've built}]
One more word for this chapter.
\end{grammarlens}

\subsection*{Your turn}
\begin{itemize}
\raggedright
  \item \emph{Say it:} \emph{khaṭvā}
  \item \emph{Say it:} \emph{khaṭvā}, once more
  \item \emph{Say it:} say \emph{khaṭvā}
  \item \emph{Recall:} say the Sanskrit for a cart, then the Sanskrit for a pot, then say \emph{khaṭvā} again
\end{itemize}

\begin{tcolorbox}[breakable,colback=teal!4,colframe=teal!35!black,title={Before you move on}]
What does \sk{खट्वा} mean? (\textquotedblleft{}A cot\textquotedblright{}.) Say it once more.
\end{tcolorbox}
